\section{Determinants}
\label{section-determinants}

\noindent
Let $R$ be a ring. Let $M$ be a finite projective $R$-module.
There exists a product decomposition $R = R_0 \times \ldots \times R_t$
such that in the corresponding decomposition $M = M_0 \times \ldots \times M_t$
of $M$ we have that $M_i$ is finite locally free of rank $i$ over $R_i$.
This follows from Algebra, Lemma \ref{algebra-lemma-finite-projective}
(to see that the rank is locally constant) and
Algebra, Lemmas \ref{algebra-lemma-disjoint-decomposition} and
\ref{algebra-lemma-disjoint-implies-product} (to decompose
$R$ into a product). In this situation we define
$$
\det(M) = \wedge^0_{R_0}(M_0) \times \ldots \times \wedge^t_{R_t}(M_t)
$$
as an $R$-module. This is a finite locally free module of rank $1$
as each term is finite locally free of rank $1$.
If $\varphi : M \to N$ is an isomorphism of finite projective
$R$-modules, then we obtain a canonical isomorphism
$$
\det(\varphi) : \det(M) \longrightarrow \det(N)
$$
of locally free modules of rank $1$. More generally, if
for all primes $\mathfrak p$ of $R$ the ranks of the free modules
$M_\mathfrak p$ and $N_\mathfrak p$ are the same, then any $R$-module
homomorphism $\varphi : M \to N$ induces an $R$-module map
$\det(\varphi) : \det(M) \to \det(N)$. Finally, if $M = N$
then $\det(\varphi) : \det(M) \to \det(M)$ is an endomorphism
of an invertible $R$-module. Since $R = \Hom_R(L, L)$
for an invertible $R$-module we may and do view $\det(\varphi)$
as an element of $R$. In this way we obtain the {\it determinant}
$$
\det : \Hom_R(M, M) \longrightarrow R
$$
which is a multiplicative map.

\begin{remark}
\label{remark-determinant-as-socle}
Let $R$ be a ring. Let $M$ be a finite projective $R$-module.
Then we can consider the graded commutative $R$-algebra
exterior algebra $\wedge^*_R(M)$ on $M$ over $R$.
A formula for $\det(M)$ is that $\det(M) \subset \wedge^*_R(M)$
is the annihilator of $M \subset \wedge^*_R(M)$.
This is sometimes useful as it does not refer to the
decomposition of $R$ into a product. Of course, to
prove this satisfies the desired properties one has to
either decompose $R$ into a product (as above), or one
has to look at the localizations at primes of $R$.
\end{remark}

\noindent
Next, we consider what happens to the determinant give
a short exact sequence of finite projective modules.

\begin{lemma}
\label{lemma-det-ses}
Let $R$ be a ring. Let
$$
0 \to M' \to M \to M'' \to 0
$$
be a short exact sequence of finite projective $R$-modules. Then there
is a canonical isomorphism
$$
\gamma : \det(M') \otimes \det(M'') \longrightarrow \det(M)
$$
\end{lemma}

\begin{proof}[First proof]
First proof. Decompose $R$ into a product of rings $R_{ij}$ such that
$M' = \prod M'_{ij}$ and $M'' = \prod M''_{ij}$ where
$M'_{ij}$ has rank $i$ and $M''_{ij}$ has rank $j$.
Of course then $M = \prod M_{ij}$ and $M_{ij}$ has rank $i + j$.
This reduces us to the case where $M'$ and $M''$ have constant
rank say $i$ and $j$. In this case we have to construct a canonical
map
$$
\wedge^i(M') \otimes \wedge^j(M'') \longrightarrow \wedge^{i + j}(M)
$$
To do this choose $m'_1, \ldots, m'_i$ in $M'$ and
$m''_1, \ldots, m''_j$ in $M''$. Denote $m_1, \ldots, m_i \in M$
the images of $m'_1, \ldots, m'_i$ and denote
$m_{i + 1}, \ldots , m_{i + j} \in M$ elements mapping to
$m''_1, \ldots, m''_j$ in $M''$. Our rule will be that
$$
m'_1 \wedge \ldots \wedge m'_i \otimes
m''_1 \wedge \ldots \wedge m''_j
\longmapsto
m_1 \wedge \ldots \wedge m_{i + j}
$$
We omit the detailed proof that this is well defined and an isomorphism.
\end{proof}

\begin{proof}[Second proof]
We will use the description of $\det(M)$, $\det(M')$, and $\det(M'')$
given in Remark \ref{remark-determinant-as-socle}. Consider the $R$-algebra
maps $\wedge^*_R(M') \to \wedge^*_R(M)$ and
$\wedge^*_R(M) \to \wedge^*_R(M'')$. The first is injective and the
second is surjective. Take an element $x' \in \det(M') \subset \wedge^*_R(M')$
and an element $x'' \in \det(M'') \subset \wedge^*_R(M'')$.
Choose an element $y'' \in \wedge^*(M)$ mapping to $x''$ and set
$$
\gamma(x' \otimes x'') = x' \wedge y'' \in \det(M) \subset \wedge^*_R(M)
$$
The reader verifies easily by looking at localizations at primes
that this well defined and an isomorphism. Moreover, this construction
gives the same map as the construction given in the first proof.
\end{proof}

\begin{lemma}
\label{lemma-det-ses-functorial}
Let $R$ be a ring. Let
$$
\xymatrix{
0 \ar[r] &
M' \ar[r] \ar[d]^u &
M \ar[r] \ar[d]^v &
M'' \ar[r] \ar[d]^w &
0 \\
0 \ar[r] &
K' \ar[r] &
K \ar[r] &
K'' \ar[r] &
0
}
$$
be a commutative diagram of finite projective $R$-modules
whose vertical arrows are isomorphisms. Then we get a commutative
diagram of isomorphisms
$$
\xymatrix{
\det(M') \otimes \det(M'') \ar[r]_-\gamma \ar[d]_{\det(u) \otimes \det(w)} &
\det(M) \ar[d]^{\det(v)} \\
\det(K') \otimes \det(K'') \ar[r]^-\gamma & \det(K)
}
$$
where the horizontal arrows are the ones constructed
in Lemma \ref{lemma-det-ses}.
\end{lemma}

\begin{proof}
Omitted. Hint: use the second construction of the maps $\gamma$
in Lemma \ref{lemma-det-ses}.
\end{proof}

\begin{lemma}
\label{lemma-det-filtration}
Let $R$ be a ring. Let
$$
K \subset L \subset M
$$
be $R$-modules such that $K$, $L/K$, and $M/L$ are finite projective
$R$-modules. Then the diagram
$$
\xymatrix{
\det(K) \otimes \det(L/K) \otimes \det(M/L) \ar[r] \ar[d] &
\det(L) \otimes \det(M/L) \ar[d] \\
\det(K) \otimes \det(M/K) \ar[r] &
\det(M)
}
$$
commutes where the maps are those of Lemma \ref{lemma-det-ses}.
\end{lemma}

\begin{proof}
Omitted. Hint: after localizing at a prime of $R$ we can assume
$K \subset L \subset M$ is isomorphic to
$R^{\oplus a} \subset R^{\oplus a + b}  \subset R^{\oplus a + b + c}$
and in this case the result is an evident computation.
\end{proof}

\begin{lemma}
\label{lemma-det-direct-sum}
Let $R$ be a ring. Let $M'$ and $M''$ be two finite projective
$R$-modules. Then the diagram
$$
\xymatrix{
\det(M') \otimes \det(M'') \ar[r]
\ar[d]_{\epsilon \cdot (\text{switch tensors})} &
\det(M' \oplus M'') \ar[d]^{\det(\text{switch summands})} \\
\det(M'') \otimes \det(M') \ar[r] &
\det(M'' \oplus M')
}
$$
commutes where $\epsilon = \det( -\text{id}_{M' \otimes M''}) \in R^*$
and the horizontal arrows are those of Lemma \ref{lemma-det-ses}.
\end{lemma}

\begin{proof}
Omitted.
\end{proof}

\begin{lemma}
\label{lemma-det-switch}
Let $R$ be a ring. Let $M$, $N$ be finite projective $R$-modules.
Let $a : M \to N$ and $b : N \to M$ be $R$-linear maps.
Then
$$
\det(\text{id} + a \circ b) = \det(\text{id} + b \circ a)
$$
as elements of $R$.
\end{lemma}

\begin{proof}
It suffices to prove the assertion after replacing $R$ by a localization
at a prime ideal. Thus we may assume $R$ is local and $M$ and $N$ are
finite free. In this case we have to prove the equality
$$
\det(I_n + AB) = \det(I_m + BA)
$$
of usual determinants of matrices where $A$ has size $n \times m$
and $B$ has size $m \times n$. This reduces to the case
of the ring $R = \mathbf{Z}[a_{ij}, b_{ji}; 1 \leq i \leq n, 1 \leq j \leq m]$
where $a_{ij}$ and $b_{ij}$ are variables and the entries of the
matrices $A$ and $B$. Taking the fraction field, this reduces to the
case of a field of characteristic zero. In characteristic zero there
is a universal polynomial expressing the determinant of a matrix
of size $\leq N$ in the traces of the powers of said matrix.
Hence it suffices to prove
$$
\text{Trace}((I_n + AB)^k) = \text{Trace}((I_m + BA)^k)
$$
for all $k \geq 1$. Expanding we see that it suffices to prove
$\text{Trace}((AB)^k) = \text{Trace}((BA)^k)$ for all $k \geq 0$.
For $k = 1$ this is the well known fact that
$\text{Trace}(AB) = \text{Trace}(BA)$. For $k > 1$
it follows from this by writing $(AB)^k = A(BA)^{k - 1}B$
and $(BA)^k = (BA)^{k - 1} A B$.
\end{proof}

\noindent
Recall that we have defined in Algebra, Section \ref{algebra-section-K-groups}
a group $K_0(R)$ as the free group on isomorphism classes
of finite projective $R$-modules modulo the relations
$[M'] + [M''] = [M' \oplus M'']$.

\begin{lemma}
\label{lemma-det}
Let $R$ be a ring. There is a map
$$
\det : K_0(R) \longrightarrow \Pic(R)
$$
which maps $[M]$ to the class of the invertible module
$\wedge^n(M)$ if $M$ is a finite locally free module of rank $n$.
\end{lemma}

\begin{proof}
This follows immediately from the constructions above and in particular
Lemma \ref{lemma-det-ses} to see that the relations are mapped to $0$.
\end{proof}





